\begin{table}[htbp]
\centering
\caption{Dynamic mechanisms at the original region center}\label{tab:r8_mechanisms}
\small
\begin{tabular}{lrrr}
\toprule
Intervention & $10^4\Delta^{\rm adj}$ & $10^4\Delta^0$ & $10^4$ relative option \\
\midrule
Baseline & 1.9674 & -2.0933 & 4.0607 \\
Tilt $a=-0.25$ & 1.5564 & -1.3521 & 2.9085 \\
Tilt $a=0.25$ & 2.3858 & -2.8501 & 5.2359 \\
Tilt $a=1$ & 3.8022 & -5.1768 & 8.9790 \\
Normalization $s=0$ & 9.1981 & 9.1981 & 0.0000 \\
Normalization $s=0.5$ & 11.6551 & 10.5243 & 1.1309 \\
Normalization $s=1.5$ & -8.9524 & -15.8240 & 6.8716 \\
Liquidation $z=0$ & -20.0244 & -24.7728 & 4.7484 \\
Liquidation $z=2$ & 21.1887 & 17.8547 & 3.3340 \\
Zero-covariance law & 4.1345 & 0.9752 & 3.1593 \\
\bottomrule
\end{tabular}
\par\vspace{3pt}\begin{minipage}{0.96\textwidth}\footnotesize Except for the named intervention, parameters are the original $(\lambda,d,k)=(0.125,0.425,2)$. Every entry reoptimizes both risk classes in both adjustment regimes. The normalization multiplier changes $b_0(u)=1/(1-u)$, not marginal consumption utility. Values shown as zero at $s=0$ are numerically negligible, not an asserted exact zero theorem. These model counterfactuals are not empirically identified causal effects.\end{minipage}
\end{table}
